\documentclass[11pt, a4paper]{article}

% ── Packages ──────────────────────────────────────────────────────────────────
\usepackage[a4paper, margin=2.5cm]{geometry}
\usepackage[utf8]{inputenc}
\usepackage[T1]{fontenc}
\usepackage{lmodern}
\usepackage{microtype}
\usepackage{parskip}
\usepackage{setspace}
\usepackage{titlesec}
\usepackage{booktabs}
\usepackage{array}
\usepackage{tabularx}
\usepackage{xcolor}
\usepackage{amsmath}
\usepackage{amssymb}
\usepackage{enumitem}
\usepackage{hyperref}
\usepackage{fancyhdr}
\usepackage{mdframed}
\usepackage{colortbl}

% ── Colours ───────────────────────────────────────────────────────────────────
\definecolor{darkblue}{HTML}{1a1a2e}
\definecolor{midblue}{HTML}{2e2e5e}
\definecolor{framebg}{HTML}{eef0ff}
\definecolor{framebd}{HTML}{aaaacc}
\definecolor{mathbg}{HTML}{f5f5ff}
\definecolor{greenbg}{HTML}{d4edda}
\definecolor{rowA}{HTML}{f0f0fa}

% ── Hyperref ──────────────────────────────────────────────────────────────────
\hypersetup{
    colorlinks=true,
    linkcolor=midblue,
    urlcolor=midblue,
    pdftitle={Clustering Results --- Technical Explanation},
}

% ── Section formatting ────────────────────────────────────────────────────────
\titleformat{\section}
  {\large\bfseries\color{darkblue}}
  {\thesection.}{0.6em}{}
  [\vspace{-0.3em}\textcolor{framebd}{\rule{\linewidth}{0.4pt}}]

\titleformat{\subsection}
  {\normalsize\bfseries\color{midblue}}
  {\thesubsection}{0.6em}{}

\titlespacing{\section}{0pt}{18pt}{6pt}
\titlespacing{\subsection}{0pt}{10pt}{4pt}

% ── Header / Footer ───────────────────────────────────────────────────────────
\pagestyle{fancy}
\fancyhf{}
\renewcommand{\headrulewidth}{0.4pt}
\fancyhead[L]{\small\textcolor{darkblue}{\textbf{Clustering Results}}}
\fancyhead[R]{\small\textcolor{darkblue}{201940185}}
\fancyfoot[C]{\small\thepage}

% ── Shaded boxes ─────────────────────────────────────────────────────────────
\newmdenv[
  backgroundcolor=mathbg,
  linecolor=framebd,
  linewidth=0.6pt,
  innerleftmargin=12pt,
  innerrightmargin=12pt,
  innertopmargin=8pt,
  innerbottommargin=8pt,
  skipabove=6pt,
  skipbelow=6pt,
]{mathbox}

\newmdenv[
  backgroundcolor=framebg,
  linecolor=framebd,
  linewidth=0.6pt,
  innerleftmargin=12pt,
  innerrightmargin=12pt,
  innertopmargin=8pt,
  innerbottommargin=8pt,
  skipabove=6pt,
  skipbelow=6pt,
]{insightbox}

% ── List style ────────────────────────────────────────────────────────────────
\setlist[itemize]{leftmargin=1.4em, itemsep=3pt, topsep=4pt}

% =============================================================================
\begin{document}
% =============================================================================

% ── Title block ───────────────────────────────────────────────────────────────
\begin{center}
  {\LARGE\bfseries\color{darkblue} Clustering Results}\\[4pt]
  {\normalsize CA Assignment 2 $\cdot$ Clustering Algorithms}\\[2pt]
  {\normalsize Algorithm: Agglomerative Hierarchical Clustering (AHC)}
\end{center}

\vspace{4pt}
\textcolor{framebd}{\rule{\linewidth}{0.8pt}}
\vspace{6pt}

% =============================================================================
\section{Clustering Algorithm: Agglomerative Hierarchical Clustering}
% =============================================================================

Agglomerative Hierarchical Clustering (AHC) is a bottom-up clustering
algorithm that builds a hierarchy of nested partitions by iteratively merging
the two closest clusters.  The process begins with $N$ singleton clusters ---
one per data point --- and terminates when all points belong to a single
cluster.  A flat partition into $K$ clusters is obtained by cutting the
resulting dendrogram at the level that produces $K$ groups.  The cut point
is selected by maximising the silhouette coefficient over $K = 1, \ldots, 10$.

\subsection{Algorithm Mechanics}

At each iteration, AHC computes the inter-cluster distance for all current
cluster pairs using a linkage criterion, then merges the pair with the
smallest distance.  The choice of linkage criterion fundamentally determines
the geometry of the resulting clusters:

\begin{itemize}
  \item \textbf{Single linkage:}
        $d(A,B) = \min_{a \in A,\, b \in B} d(a,b)$.
        Prone to chaining; produces elongated, non-compact clusters.

  \item \textbf{Average linkage:}
        $d(A,B) = \frac{1}{|A||B|}\sum_{a \in A}\sum_{b \in B} d(a,b)$.
        Less prone to chaining than single linkage but still susceptible on
        text data where one domain greatly outnumbers others.

  \item \textbf{Ward linkage (used in this work):}  merges the pair $(A, B)$
        that minimises the increase in total within-cluster sum of squares.
        Produces compact, balanced, well-separated clusters.
\end{itemize}

\subsection{Why AHC?}

\begin{itemize}
  \item \textbf{No assumption about cluster shape.}  Unlike $K$-means
        (MacQueen, 1967), which assumes spherical, equally-sized clusters,
        AHC can recover clusters of arbitrary shape and size.  The dataset
        exhibits strong size imbalance: Cluster~2 (Shakespeare, 766 documents)
        and Cluster~3 (Wikipedia, 89 documents).

  \item \textbf{Deterministic output.}  AHC produces the same result on
        every run, unlike $K$-means which depends on random centroid
        initialisation.  This is important for reproducibility.

  \item \textbf{No need to pre-specify $K$.}  AHC builds the complete
        dendrogram in a single pass; the optimal $K$ is selected post-hoc by
        silhouette analysis across $K = 1$ to $10$.

  \item \textbf{Literature support.}  AHC is a standard and well-validated
        algorithm for text clustering with dense embeddings.  Petukhova et
        al.\ (2024) include AHC as one of their primary benchmark algorithms
        and report competitive performance with BERT-family embeddings.
\end{itemize}

% =============================================================================
\section{Configuration: Ward Linkage and Euclidean Distance}
% =============================================================================

\subsection{Ward Linkage}

Ward linkage (Ward, 1963) selects the merge that minimises the increase in total
within-cluster sum of squares (WCSS).  Formally, the cost of merging clusters
$A$ and $B$ is:

\begin{mathbox}
\[
  \Delta(A,\,B)
  \;=\;
  \frac{|A|\,|B|}{|A|+|B|}
  \;\|\boldsymbol{\mu}_A - \boldsymbol{\mu}_B\|^{2}
\]
\end{mathbox}

where $\boldsymbol{\mu}_A$ and $\boldsymbol{\mu}_B$ are the cluster centroids
and $|A|$, $|B|$ are cluster sizes.  Ward is the only linkage criterion that
directly optimises cluster compactness at every merge step.  This produces
balanced, well-separated clusters and is robust against the
\textbf{chaining problem} --- the tendency of average and single linkage to
greedily merge individual points into one large elongated cluster rather than
forming distinct groups.

On this dataset, empirical investigation confirmed the chaining problem: average
linkage consistently collapsed the clustering to $K=2$ regardless of metric,
absorbing the Wikipedia sub-corpus into the larger modern English cluster and
ignoring its thematic distinctiveness.  Ward linkage recovered the semantically
correct $K=3$ solution.

\subsection{Euclidean Distance}

Ward linkage is mathematically defined \emph{only} for Euclidean distance ---
the formula for $\Delta(A,B)$ above uses squared Euclidean norms between
centroids.  Applying Ward with any other metric (e.g., cosine) is not
mathematically valid and produces undefined behaviour.

The use of Euclidean distance is appropriate here for two reasons.  First,
L2 normalisation is applied to the raw embeddings before PCA.  For unit vectors
in $\mathbb{R}^{384}$, Euclidean and cosine distances are monotonically
equivalent:

\begin{mathbox}
\[
  d_{\mathrm{Euclidean}}(\hat{\mathbf{a}},\,\hat{\mathbf{b}})^{2}
  = 2\bigl(1 - \cos(\hat{\mathbf{a}},\,\hat{\mathbf{b}})\bigr)
\]
\end{mathbox}

PCA of unit vectors preserves this angular geometry in the principal component
subspace, so Euclidean distance in PCA space remains semantically meaningful.
Second, the PCA output is deliberately left un-normalised, so Euclidean
distances in the 50-dimensional PCA space weight the earlier, higher-variance
principal components more heavily --- since these components capture the most
semantically significant axes of variation, this variance weighting is a useful
inductive bias for cluster discovery.

% =============================================================================
\section{Silhouette Analysis}
% =============================================================================

The silhouette coefficient is an \textbf{internal} cluster validity index that
quantifies how well each data point fits its assigned cluster relative to the
nearest alternative cluster (Rousseeuw, 1987).  For data point $i$:

\begin{mathbox}
\begin{align*}
  a(i) &= \text{mean intra-cluster Euclidean distance from } i
           \text{ to all other points in its cluster} \\[4pt]
  b(i) &= \text{mean Euclidean distance from } i \text{ to all points
           in the nearest other cluster} \\[4pt]
  s(i) &= \frac{b(i) - a(i)}{\max\bigl(a(i),\,b(i)\bigr)}
           \;\in\; [-1,\, +1]
\end{align*}
\end{mathbox}

The overall silhouette score is:

\begin{mathbox}
\[
  S \;=\; \frac{1}{N}\sum_{i=1}^{N} s(i)
\]
\end{mathbox}

A score near $+1$ indicates well-separated, compact clusters; near $0$
indicates overlapping clusters; negative values indicate misassignment.  Both
$a(i)$ and $b(i)$ are computed using the same Euclidean metric used for
clustering, ensuring a fully consistent and internally valid evaluation.

\subsection{Results for $K = 1$ to $10$}

\begin{table}[h!]
\centering
\renewcommand{\arraystretch}{1.35}
\begin{tabular}{>{\centering\arraybackslash}p{2cm}
                >{\centering\arraybackslash}p{5cm}
                p{7.5cm}}
  \toprule
  \textbf{$K$} & \textbf{Silhouette Coefficient} & \textbf{Notes} \\
  \midrule
  1  & N/A      & Undefined --- requires at least 2 clusters \\
  2  & 0.108900 & \\
  \rowcolor{greenbg}
  3  & \textbf{0.116243} & \textbf{Global optimum (peak)} \\
  4  & 0.087409 & \\
  5  & 0.064581 & \\
  6  & 0.055980 & \\
  7  & 0.056356 & \\
  8  & 0.064259 & \\
  9  & 0.070488 & \\
  10 & 0.077989 & \\
  \bottomrule
\end{tabular}
\caption{Silhouette coefficients for AHC (Euclidean, Ward) over $K = 1$
         to $10$.  $K=3$ is highlighted as the global optimum.}
\end{table}

% =============================================================================
\section{Optimal Value of $K$ and Justification}
% =============================================================================

\subsection{Quantitative Justification: Silhouette Analysis}

The silhouette analysis unambiguously identifies $\mathbf{K=3}$ as the optimal
number of clusters, with a peak coefficient of $\mathbf{0.116243}$.  The score
at $K=3$ exceeds that at $K=2$ (0.108900) and all higher values of $K$, which
decrease monotonically to $0.055980$ at $K=6$ before slightly recovering at
$K=9$--$10$.  The shape of the curve --- a clear peak at $K=3$ followed by
sustained decline --- is a strong quantitative signal that the data contains
three natural groupings at the granularity captured by the MiniLM embedding and
the 50-dimensional PCA feature space.

\subsection{Qualitative Justification: Semantic Cluster Interpretation}

The three clusters produced at $K=3$ have clear, interpretable semantic
identities:

\begin{table}[h!]
\centering
\renewcommand{\arraystretch}{1.5}
\begin{tabularx}{\textwidth}{>{\centering\bfseries}p{1.5cm}
                              >{\centering}p{2cm}
                              >{\centering}p{3cm}
                              X}
  \toprule
  \textbf{Cluster} & \textbf{Size} & \textbf{Domain} &
  \textbf{Linguistic Characteristics} \\
  \midrule
  1 & 878 (50.6\%) & Reviews + News &
    Modern English prose: restaurant and business reviews (Yelp-style) and
    CNN-style news articles.  Informal to formal register, contemporary
    vocabulary, first and third person narration. \\[4pt]
  2 & 766 (44.2\%) & Shakespeare (Coriolanus) &
    Early Modern English: play dialogue from Coriolanus.  Archaic vocabulary,
    poetic metre, character prefixes (e.g., \textsc{Coriolanus:},
    \textsc{Aufidius:}), iambic pentameter constructions. \\[4pt]
  3 & 89\ \ (5.1\%)  & Wikipedia (Notre Dame) &
    Encyclopaedic prose: factual Wikipedia-style articles about Notre Dame
    University.  Highly formal, third-person, dense noun phrases,
    institutional and academic vocabulary. \\
  \bottomrule
\end{tabularx}
\caption{Semantic interpretation of the three clusters produced by AHC at
         $K=3$.}
\end{table}

The merging of reviews and news into a single cluster (Cluster~1) is
semantically well-founded: both sub-corpora are written in contemporary English,
use similar vocabulary (named entities, everyday objects, locations), and share
syntactic patterns characteristic of modern prose.  In MiniLM's embedding space,
their angular separation is smaller than their collective angular separation from
Shakespeare.  This is not a failure of the algorithm but an \textbf{accurate
reflection of the linguistic properties of the data}.

If a four-cluster solution were desired for theoretical reasons, $K=4$ would be
the next candidate; however, its silhouette score of $0.087409$ --- a 25\% drop
from $K=3$ --- provides no quantitative support for the additional split.  The
silhouette criterion is thus unambiguous: $K=3$ is optimal.

\subsection{Contextualisation Against the Literature}

Silhouette scores for text clustering with BERT-family embeddings are
structurally low compared to geometric or image datasets, because sentence
embeddings reside in a dense, high-dimensional space where cluster boundaries
are inherently fuzzy.  Petukhova et al.\ (2024) report silhouette scores for
BERT-family models ranging from $0.016$ to $0.224$ across five diverse text
datasets.  The score of $0.116243$ achieved here falls comfortably within this
range --- comparable to BERT + AHC on the SyskillWebert dataset
($\mathrm{SS} = 0.152$) and notably higher than BERT + $k$-means on the
20Newsgroups dataset ($\mathrm{SS} = 0.048$).

This confirms that the result is realistic, well-calibrated, and consistent
with state-of-the-art benchmarks.  A silhouette score approaching $1.0$ would
in fact be cause for suspicion: such values typically indicate that a secondary
non-linear dimensionality reduction step such as UMAP (McInnes et al., 2018)
with $\mathtt{min\_dist}=0$ is artificially inflating the metric by compressing
points into geometrically tight blobs that are geometrically but not semantically
well-separated.

% =============================================================================
\section{Additional Considerations and Limitations}
% =============================================================================

\subsection{Row Count Discrepancy}

The dataset specification states 1{,}760 instances; however, only 1{,}733 rows
are successfully loaded.  The discrepancy arises because a subset of rows
contains literal newline characters embedded within the \texttt{Sentence} field.
When \texttt{pandas} reads the tab-separated file, these embedded newlines are
interpreted as line breaks, producing malformed records that are silently
discarded via \texttt{on\_bad\_lines=`skip'}.  The \texttt{label.txt} output
contains exactly 1{,}733 labels, one per successfully loaded instance, in
original row order.

\subsection{Cluster Size Imbalance}

Cluster~3 (Wikipedia) contains only 89 documents ($\approx 5.1\%$ of the
dataset).  This imbalance is a genuine property of the data rather than an
algorithmic artefact.  Ward linkage is the linkage criterion most resistant to
distortion by cluster size imbalance: its variance-minimisation objective does
not favour equal-sized clusters.  The small Wikipedia cluster is recovered
because its embedding centroid is sufficiently distant from both the Shakespeare
and reviews/news centroids in PCA space to justify its existence as a separate
group under the silhouette criterion.

\subsection{Embedding Model Limitations}

MiniLM-L12-v2 produces 384-dimensional embeddings, compared to 768 dimensions
for larger models such as \texttt{all-mpnet-base-v2}.  Larger embeddings capture
finer-grained semantic distinctions and may improve the separability of the
reviews and news sub-corpora within Cluster~1.  However, the assignment
constraint of no cloud API dependency and the practical requirement for local
inference make MiniLM the appropriate choice.  Petukhova et al.\ (2024) confirm
that BERT-family models represent the best open-source option for text
clustering, outperforming TF-IDF baselines and performing comparably to
much larger open-source LLMs.

\vspace{14pt}
\noindent\textcolor{midblue}{\small\textbf{References:}\\[3pt]
MacQueen, J.\ (1967).\ Some methods for classification and analysis of
multivariate observations. \textit{Proceedings of the 5th Berkeley Symposium
on Mathematical Statistics and Probability}, 1, 281--297.\\[3pt]
McInnes, L., Healy, J., \& Melville, J.\ (2018).\ UMAP: Uniform manifold
approximation and projection for dimension reduction.
\textit{arXiv:1802.03426}.\\[3pt]
Petukhova, A., Matos-Carvalho, J.\,P., \& Fachada, N.\ (2024).\ Text clustering
with LLM embeddings. \textit{arXiv:2403.15112}.\\[3pt]
Reimers, N., \& Gurevych, I.\ (2019).\ Sentence-BERT: Sentence embeddings using
Siamese BERT-networks. \textit{Proceedings of EMNLP 2019}, 3982--3992.\\[3pt]
Rousseeuw, P.\,J.\ (1987).\ Silhouettes: A graphical aid to the interpretation
and validation of cluster analysis. \textit{Journal of Computational and Applied
Mathematics}, 20, 53--65.\\[3pt]
Ward, J.\,H.\ (1963).\ Hierarchical grouping to optimize an objective function.
\textit{Journal of the American Statistical Association}, 58(301), 236--244.}

\end{document}